\documentclass[10pt,a4paper,twocolumn]{article}

\usepackage[utf8]{inputenc}
\usepackage[T1]{fontenc}
\usepackage{lmodern}
\usepackage[margin=2cm, columnsep=0.8cm]{geometry}
\usepackage{amsmath, amssymb}
\usepackage{graphicx}
\usepackage{booktabs}
\usepackage[numbers,sort&compress]{natbib}
\usepackage[hidelinks]{hyperref}

\title{\bfseries A Two-Column Research Paper}
\author{First Author\thanks{Multimedia University, Malaysia} \and Second Author\footnotemark[1]}
\date{}

\begin{document}

\twocolumn[
  \begin{@twocolumnfalse}
    \maketitle
    \begin{abstract}
      Write a concise summary of the paper here. The abstract spans both columns so it stands out from the body text.
    \end{abstract}
    \vspace{1em}
  \end{@twocolumnfalse}
]

\section{Introduction}
Two-column layouts are common for conference papers. Text flows from the left column to the right column automatically \cite{ref1}.

\section{Background}
Explain the necessary background here. Wide figures can span both columns with \texttt{figure*}.

\section{Method}
\begin{equation}
    \mathcal{L}(\theta) = -\frac{1}{N} \sum_{i=1}^{N} \log p_\theta(y_i \mid x_i)
\end{equation}

\section{Results}
\begin{table}[h]
    \centering
    \caption{Results summary.}
    \begin{tabular}{lr}
        \toprule
        Setting & Score \\
        \midrule
        A & 71.2 \\
        B & 78.9 \\
        \bottomrule
    \end{tabular}
\end{table}

\section{Conclusion}
Conclude the paper and outline future work.

\bibliographystyle{unsrtnat}
\bibliography{references}

\end{document}
